\documentclass[11pt,a4paper]{article}

\usepackage[margin=1in]{geometry}
\usepackage{amsmath,amssymb,amsthm}
\usepackage{booktabs}
\usepackage{hyperref}

\newcommand{\ebasis}[1]{\mathbf{e}_{#1}}
\newcommand{\conj}[1]{\overline{#1}}
\newcommand{\dOmega}{\,\mathrm{d}\Omega}
% \eth is provided by amssymb

\title{Canonical Components and Generalized Spherical Harmonics\\
\large Conventions, symmetries and admissible operations}
\author{Notes for the GSHTrans field algebra}
\date{}

\begin{document}
\maketitle

\begin{abstract}
This note fixes the conventions used by the GSHTrans field algebra and
derives the facts that the type system is intended to enforce: how the
generalized spherical harmonic upper index is carried by a tensor
component, which algebraic operations on components are covariant, and
how reality and permutation symmetry reduce the set of components that
must be stored.  Every convention-dependent statement is flagged; these
are the points that need checking against Phinney \& Burridge before any
code is written.
\end{abstract}

\tableofcontents

\section{The canonical basis}
\label{sec:basis}

At each point of the sphere let $\hat{\mathbf{r}},
\hat{\boldsymbol{\theta}}, \hat{\boldsymbol{\phi}}$ be the usual
orthonormal spherical polar frame.  The canonical basis is
\begin{equation}
  \ebasis{\pm} = \mp \tfrac{1}{\sqrt{2}}
    \left( \hat{\boldsymbol{\theta}} \pm i \hat{\boldsymbol{\phi}} \right),
  \qquad
  \ebasis{0} = \hat{\mathbf{r}} .
  \label{eq:basis}
\end{equation}
The index $\alpha \in \{-1, 0, +1\}$, written $\{-,0,+\}$ where no
confusion arises.

\paragraph{Convention check.}  The sign $\mp$ in \eqref{eq:basis} is the
first thing to verify.  Some authors use $\ebasis{\pm} = \pm(\hat{\boldsymbol{\theta}}
\pm i\hat{\boldsymbol{\phi}})/\sqrt{2}$, which flips the sign of every
odd-rank reality phase below.

Two properties follow immediately and are used throughout.  First, complex
conjugation of the basis:
\begin{equation}
  \conj{\ebasis{\alpha}} = (-1)^{\alpha}\, \ebasis{-\alpha},
  \label{eq:basisconj}
\end{equation}
which reproduces $\conj{\ebasis{\pm}} = -\ebasis{\mp}$ and
$\conj{\ebasis{0}} = \ebasis{0}$.  Second, the metric
$g_{\alpha\beta} = \ebasis{\alpha} \cdot \ebasis{\beta}$:
\begin{equation}
  g_{\alpha\beta} = (-1)^{\alpha}\, \delta_{\alpha+\beta,\,0},
  \qquad\text{that is}\qquad
  g_{+-} = g_{-+} = -1, \quad g_{00} = 1,
  \label{eq:metric}
\end{equation}
all other entries vanishing.  Note $g$ is its own inverse.  We work
throughout with \emph{contravariant} components only; the metric
\eqref{eq:metric} is applied explicitly wherever a contraction is taken,
and index position is not tracked.

\section{Canonical components of a tensor}
\label{sec:components}

A rank-$p$ tensor field on the sphere is expanded as
\begin{equation}
  \mathbf{T}(\theta,\phi) = \sum_{\alpha_1 \ldots \alpha_p}
    T^{\alpha_1 \ldots \alpha_p}(\theta,\phi)\,
    \ebasis{\alpha_1} \otimes \cdots \otimes \ebasis{\alpha_p} ,
\end{equation}
so a rank-$p$ tensor has $3^p$ scalar components, each labelled by a
\emph{multi-index} $(\alpha_1 \ldots \alpha_p) \in \{-1,0,1\}^p$.

Under a rotation of the tangent frame through an angle $\psi$ about
$\hat{\mathbf{r}}$ the basis vectors transform as $\ebasis{\pm} \mapsto
e^{\mp i\psi} \ebasis{\pm}$, $\ebasis{0} \mapsto \ebasis{0}$.  Hence the
component $T^{\alpha_1 \ldots \alpha_p}$ acquires the phase
$e^{-iN\psi}$ with
\begin{equation}
  \boxed{\;N = \sum_{i=1}^{p} \alpha_i \;}
  \label{eq:N}
\end{equation}
and it is this $N$, not the multi-index itself, that is the upper index
of the generalized spherical harmonic expansion of that component
(Section~\ref{sec:gsh}).

Two consequences matter for the design of the library.

\begin{enumerate}
  \item For a vector ($p=1$) the multi-index and the upper index
    coincide, $N = \alpha$.  For $p \geq 2$ they do not: distinct
    components share an upper index.  A rank-2 tensor has three
    components with $N=0$, namely $(-+)$, $(00)$ and $(+-)$.  A
    collection of components labelled only by $N$ therefore does
    \emph{not} determine the tensor.
  \item The number of components of a rank-$p$ tensor with a given upper
    index $N$ is the trinomial coefficient
    $\binom{p}{N}_{\!2}$; for $p=2$ these are $1,2,3,2,1$ for
    $N=-2,\ldots,+2$.
\end{enumerate}

\begin{table}[h]
  \centering
  \begin{tabular}{ccc}
    \toprule
    $N$ & components & count \\
    \midrule
    $-2$ & $(--)$ & 1 \\
    $-1$ & $(-0),\ (0-)$ & 2 \\
    $ 0$ & $(-+),\ (00),\ (+-)$ & 3 \\
    $+1$ & $(0+),\ (+0)$ & 2 \\
    $+2$ & $(++)$ & 1 \\
    \bottomrule
  \end{tabular}
  \caption{Upper index of each component of a rank-2 tensor.}
  \label{tab:rank2}
\end{table}

\paragraph{Tensor product.}  If $S$ has multi-index
$(\alpha_1 \ldots \alpha_p)$ with upper index $N_S$ and $T$ has
$(\beta_1 \ldots \beta_q)$ with $N_T$, then the component of
$S \otimes T$ with the concatenated multi-index has upper index
$N_S + N_T$ by \eqref{eq:N}.  The rule ``upper indices add under
pointwise multiplication'' is exactly this statement.

\section{Reality}
\label{sec:reality}

Suppose $\mathbf{T}$ is a real tensor field.  Writing
$\mathbf{T} = \conj{\mathbf{T}}$ and using \eqref{eq:basisconj} on each
factor,
\begin{align}
  \sum_{\alpha_1 \ldots \alpha_p} T^{\alpha_1 \ldots \alpha_p}\,
    \ebasis{\alpha_1} \otimes \cdots
  &= \sum_{\alpha_1 \ldots \alpha_p} \conj{T^{\alpha_1 \ldots \alpha_p}}\,
     (-1)^{\alpha_1 + \cdots + \alpha_p}\,
     \ebasis{-\alpha_1} \otimes \cdots ,
\end{align}
and relabelling $\alpha_i \to -\alpha_i$ on the right gives the
\textbf{reality condition}
\begin{equation}
  \boxed{\;
  T^{-\alpha_1 \ldots -\alpha_p}
    = (-1)^{N}\, \conj{T^{\alpha_1 \ldots \alpha_p}},
  \qquad N = \textstyle\sum_i \alpha_i .
  \;}
  \label{eq:reality}
\end{equation}
For a vector this reads $u^{0} = \conj{u^{0}}$ and
$u^{-} = -\,\conj{u^{+}}$: the radial component is a real field, and the
two transverse components are not independent.

\subsection{Which components must be stored}
\label{sec:orbits}

Equation \eqref{eq:reality} pairs each multi-index with its
\emph{negation} $(\alpha_i) \mapsto (-\alpha_i)$.  The components that
must be stored are one representative of each orbit of this involution.

\begin{itemize}
  \item The only fixed point of the negation map is the all-zero
    multi-index, which has $N = 0$ and, by \eqref{eq:reality}, satisfies
    $T^{0\ldots 0} = \conj{T^{0 \ldots 0}}$: it is a \emph{real-valued}
    field.
  \item Every other multi-index lies in an orbit of size two; one member
    is stored as a complex field and the other is recovered by
    \eqref{eq:reality}.
\end{itemize}

The storage cost in real numbers is therefore
$1 + 2 \cdot \tfrac{1}{2}(3^p - 1) = 3^p$, which is precisely the number
of real degrees of freedom of a real rank-$p$ tensor in three
dimensions.  The reduction loses nothing and saves a factor of two
against storing all $3^p$ components as complex fields.

\begin{table}[h]
  \centering
  \begin{tabular}{lrrrr}
    \toprule
    rank $p$ & components $3^p$ & stored fields & of which real & reals stored \\
    \midrule
    0 & 1  & 1  & 1 & 1 \\
    1 & 3  & 2  & 1 & 3 \\
    2 & 9  & 5  & 1 & 9 \\
    4 & 81 & 41 & 1 & 81 \\
    \bottomrule
  \end{tabular}
  \caption{Storage for a real tensor field, per grid point, under the
    orbit rule of Section~\ref{sec:orbits}.  Complex tensors of the same
    rank require $2 \cdot 3^p$ reals.}
  \label{tab:storage}
\end{table}

\paragraph{Why ``$N \geq 0$'' is not the rule.}  For a vector the orbits
are $\{0\}$ and $\{+,-\}$, and selecting components with $N \geq 0$ picks
exactly one from each.  This coincidence fails at rank 2: $(+-)$ and
$(-+)$ both have $N = 0$ and form a negation pair, so the criterion
$N \geq 0$ would retain both.  The selection rule must be stated on
multi-indices, not on upper indices.

\subsection{Interaction with permutation symmetry}
\label{sec:symmetry}

If the tensor is symmetric (or antisymmetric) under some subgroup $P$ of
index permutations, the stored set is one representative per orbit of the
group generated by $P$ together with the negation map.  A component whose
$P$-orbit contains its own negation is constrained by \eqref{eq:reality}
to be real, or purely imaginary if the permutation relating them carries
a minus sign.

Such self-paired components always have $N = 0$: permutation preserves
the multiset sum while negation changes its sign, so a component fixed by
the combined operation satisfies $N = -N$.  Since $(-1)^{0} = +1$, a
self-paired component of a \emph{symmetric} tensor is always real; purely
imaginary components arise only in the antisymmetric case.

\paragraph{Worked check: symmetric real rank 2.}  The six independent
symmetric components are $(--), (-0), (-+), (00), (0+), (++)$.  Under the
combined group,
\begin{center}
  \begin{tabular}{lll}
    \toprule
    orbit & size & stored as \\
    \midrule
    $\{(00)\}$                    & 1 & real field \\
    $\{(-+), (+-)\}$              & 2 & real field (self-paired, $N=0$) \\
    $\{(--), (++)\}$              & 2 & complex field \\
    $\{(-0), (0-), (+0), (0+)\}$  & 4 & complex field \\
    \bottomrule
  \end{tabular}
\end{center}
giving $1 + 1 + 2 + 2 = 6$ real numbers per point, which is the number of
independent entries of a real symmetric $3\times3$ matrix.  That the
orbit construction reproduces this without special-casing is a useful
check that the abstraction is the right one.

\section{Generalized spherical harmonics}
\label{sec:gsh}

A field $f$ of upper index $N$ is expanded as
\begin{equation}
  f(\theta,\phi) = \sum_{l \geq |N|} \sum_{m=-l}^{l}
    f^{N}_{lm}\, Y^{N}_{lm}(\theta,\phi),
  \qquad
  f^{N}_{lm} = \int_{S^2} \conj{Y^{N}_{lm}}\, f \dOmega ,
  \label{eq:expansion}
\end{equation}
with
\begin{equation}
  Y^{N}_{lm}(\theta,\phi)
    = \sqrt{\frac{2l+1}{4\pi}}\; d^{l}_{Nm}(\theta)\, e^{im\phi}
  \label{eq:gsh}
\end{equation}
in the \texttt{Ortho} normalisation, where $d^{l}_{Nm}$ is the real
Wigner $d$-function.  These are orthonormal on the sphere,
$\int \conj{Y^{N}_{lm}} Y^{N}_{l'm'} \dOmega = \delta_{ll'}\delta_{mm'}$.
The sum over $l$ starts at $|N|$ because $d^{l}_{Nm}$ vanishes
identically for $l < |N|$.

\paragraph{Convention check --- now settled.}  This equation carried
$d^{l}_{mN}$, with the indices the other way round, and a note that the
order should be confirmed against \texttt{src/Wigner.h}.  It has been, and
the order above is the corrected one: the stored value is
$\sqrt{(2l+1)/4\pi}\,d^{l}_{Nm}$, which is exactly the scalar harmonic
$X^{N}_{lm}$ of Dahlen \& Tromp (C.117), built on their generalized
Legendre function $P^{N}_{lm}$ with the \emph{upper index first}.  The
check discriminates, because $d^{l}_{mN} = (-1)^{m-N} d^{l}_{Nm}$ differ
by a sign whenever $m-N$ is odd.  Woodhouse \& Deuss [41] use the same
index order without the normalising factor.  The
normalisation in \eqref{eq:gsh} was already pinned by the $l=0$ behaviour
of the existing transform: taking $f$ constant, \eqref{eq:expansion}
gives $f^{0}_{00} = 2\sqrt{\pi} f$ and $f = f^{0}_{00}/(2\sqrt{\pi})$,
which is what \texttt{GaussLegendreGrid} implements on its one-point
grid.

\subsection{Reality in the spectral domain}
\label{sec:spectralreality}

Using $d^{l}_{-N,-m} = (-1)^{m-N} d^{l}_{Nm}$, which is Dahlen \& Tromp
(C.118), and the reality of $d^{l}_{Nm}$, equation \eqref{eq:gsh} gives
\begin{equation}
  \conj{Y^{N}_{lm}} = (-1)^{m+N}\, Y^{-N}_{l,-m} .
  \label{eq:Yconj}
\end{equation}
Two distinct consequences follow, and it is worth keeping them apart.

\paragraph{Basis-level.}  For any field $f$ of upper index $N$, the
coefficients of $\conj{f}$, which has upper index $-N$, are
\begin{equation}
  \left(\conj{f}\right)^{-N}_{l,-m} = (-1)^{m-N}\, \conj{f^{N}_{lm}} .
  \label{eq:basiclevel}
\end{equation}
This is the relation exercised by \texttt{TestRealFieldSymmetry}.

\paragraph{Component-level.}  For a component pair of a \emph{real}
tensor, \eqref{eq:reality} supplies an extra factor $(-1)^{N}$, and the
two combine to give
\begin{equation}
  \boxed{\;
    T^{-N}_{l,-m} = (-1)^{m}\, \conj{T^{N}_{lm}}
  \;}
  \label{eq:complevel}
\end{equation}
for the coefficients of the negation-paired components.  At $N=0$ this
reduces to the familiar condition $f_{l,-m} = (-1)^m \conj{f_{lm}}$ for a
real scalar field, whose reduced $m \geq 0$ storage the library already
implements.

\section{Admissible algebraic operations}
\label{sec:algebra}

An operation on canonical components is admissible only if its result is
again a component of definite upper index, that is, only if it commutes
with the frame rotation $\ebasis{\pm} \mapsto e^{\mp i \psi}
\ebasis{\pm}$.  Writing $[f] = N$ for the upper index of $f$:

\begin{center}
\begin{tabular}{lll}
  \toprule
  operation & upper index of result & restriction \\
  \midrule
  $f + g$, $f - g$      & $N$            & requires $[f] = [g] = N$ \\
  $f g$                 & $[f] + [g]$    & --- \\
  $f / g$               & $[f]$          & requires $[g] = 0$ \\
  $s f$, $s \in \mathbb{C}$ & $[f]$      & --- \\
  $\conj{f}$            & $-[f]$         & --- \\
  $|f|$, $|f|^2$        & $0$            & --- \\
  $\operatorname{Re} f$, $\operatorname{Im} f$ & $0$ & requires $[f] = 0$ \\
  $F(f)$, $F$ nonlinear & $0$            & requires $[f] = 0$ \\
  $\int_{S^2} f \dOmega$ & ---           & requires $[f] = 0$ \\
  \bottomrule
\end{tabular}
\end{center}

Three of these deserve comment, because the present implementation
violates them.

\begin{itemize}
  \item \textbf{Conjugation reverses the upper index.}  A quantity
    carrying $e^{-iN\psi}$ has a conjugate carrying $e^{+iN\psi}$.  This
    is consistent with \eqref{eq:reality}: $u^{-}$, of upper index $-1$,
    is $-\conj{u^{+}}$ and $u^{+}$ has upper index $+1$.

  \item \textbf{Real and imaginary parts are defined only at $N = 0$.}
    For $N \neq 0$ the split of $f$ into real and imaginary parts is not
    preserved by a frame rotation, so neither part is a component of any
    tensor.  What \emph{is} available at any $N$ is the modulus, which is
    a genuine scalar.

  \item \textbf{Nonlinear functions likewise require $N = 0$.}  Only for
    $N=0$, where the component is invariant, is $F(f)$ again a
    component.  Integration over the sphere is subject to the same
    restriction, since $\int f \dOmega$ vanishes identically unless
    $N = 0$.
\end{itemize}

\paragraph{Contraction.}  Contracting the $j$th and $k$th slots of a
rank-$p$ tensor uses the metric \eqref{eq:metric}:
\begin{equation}
  \left(\operatorname{tr}_{jk} T\right)^{\ldots}
   = \sum_{\alpha\beta} g_{\alpha\beta}\,
      T^{\ldots\, \alpha \,\ldots\, \beta \,\ldots}
   = \sum_{\alpha} (-1)^{\alpha}\,
      T^{\ldots\, \alpha \,\ldots\, -\alpha \,\ldots} ,
\end{equation}
the contracted pair contributing $\alpha + (-\alpha) = 0$ to the upper
index, so the result has the upper index of the surviving slots, as it
must.

\section{Raising and lowering operators}
\label{sec:eth}

These are not required for the first stage of the field algebra but are
recorded here because they constrain its design: they are the operations
that connect different upper indices, and they are diagonal in $(l,m)$
rather than local in $(\theta,\phi)$.

The operators $\eth$ and $\conj{\eth}$ raise and lower the upper index by
one.  Acting on a basis function of upper index $N$,
\begin{equation}
  \eth\, Y^{N}_{lm} = -\sqrt{(l-N)(l+N+1)}\; Y^{N+1}_{lm},
  \qquad
  \conj{\eth}\, Y^{N}_{lm} = +\sqrt{(l+N)(l-N+1)}\; Y^{N-1}_{lm},
  \label{eq:eth}
\end{equation}
so that in the spectral domain each is a multiplication of the
coefficients by an $l$-dependent factor, with no coupling between
different $(l,m)$.

\paragraph{Convention check.}  The overall signs in \eqref{eq:eth}, and
the relation between $Y^{N}_{lm}$ and the spin-weighted harmonics
${}_{s}Y_{lm}$ (whether $s = N$ or $s = -N$), vary between sources and
must be fixed against Phinney \& Burridge.  The magnitudes are not in
doubt.

The surface gradient of a scalar field $f$ has canonical components
proportional to $\eth f$ and $\conj{\eth} f$, of upper index $+1$ and
$-1$ respectively, with vanishing radial component; surface divergence
and curl of a vector field are the corresponding contractions.  The
practical consequence for the library is that an expression such as ``the
gradient of a product'' is necessarily evaluated in two representations:
the product is local in $(\theta,\phi)$, the gradient is local in
$(l,m)$.

\section{Summary of facts the type system should enforce}

\begin{enumerate}
  \item A component is labelled by a multi-index; its upper index is the
    signed sum \eqref{eq:N} and is a compile-time quantity.
  \item Addition requires equal upper indices; multiplication adds them;
    division requires a denominator of upper index zero.
  \item Conjugation maps $N \mapsto -N$.
  \item Real part, imaginary part, nonlinear functions and integration
    over the sphere require $N = 0$.
  \item A real-valued component field can occur only at $N = 0$.
  \item For a real tensor, the stored components are one per orbit of
    index negation, combined with any permutation symmetry; the
    all-zero component is real-valued and the remainder are complex.
\end{enumerate}

\end{document}
